% Глава: 03-shell
% Статус: наполнено (первый прогон).
\chapter{Командная оболочка}
\label{ch:03-shell}

\section{Что такое оболочка}

\textbf{Оболочка} (shell) --- это программа-интерпретатор, которая читает
введённые вами команды, разбирает их и запускает соответствующие программы.
Наиболее распространённая оболочка Linux --- \cmd{bash}
(Bourne Again Shell). Оболочка --- это \emph{мост} между пользователем
и ядром операционной системы: она переводит ваш набор символов в вызовы
системных служб (\emph{syscalls}) и управляет жизненным циклом
запущенных программ.

\begin{sloppypar}
Все команды этой главы выполняются на вашей ВМ на сервере
\fpath{pglab.nsu.ru}.
\end{sloppypar}

\begin{figure}[htb]
  \centering
  \begin{tikzpicture}[node distance=1.6cm]
    \node[ext] (user) {Пользователь};
    \node[process, right=of user] (shell) {Оболочка\\[0.2em] bash};
    \node[process, right=of shell] (proc) {Запущенная\\[0.2em] программа};
    \node[process, below=of shell] (kernel) {Ядро Linux\\[0.2em]
        (syscalls)};
    \draw[flow, <->] (user) -- node[above, font=\small]{текст команд} (shell);
    \draw[flow, <->] (shell) -- node[right, font=\small]{fork/exec} (proc);
    \draw[flow, <->] (shell) -- node[left, font=\small]{open / read / write / exit} (kernel);
    \draw[flow, <->] (proc) -- (kernel);
  \end{tikzpicture}
  \caption{Оболочка парсит команды, создаёт подпроцессы
    (\emph{fork/exec}) и через ядро управляет ресурсами.}
  \label{fig:shell-kernel}
\end{figure}

Важно разделять \textbf{встроенные команды} (встроенные в саму оболочку:
\cmd{cd}, \cmd{echo}, \cmd{export}, \cmd{exit}) и \textbf{внешние
программы} (отдельные исполняемые файлы в \fpath{/usr/bin}, например
\cmd{ls}, \cmd{grep}). Встроенные команды работают быстрее и могут
менять состояние самой сессии (например, \cmd{cd} меняет каталог именно
\emph{текущей} оболочки).

\section{Сессии и подпроцессы}

Каждое окно терминала --- это отдельная \textbf{сессия} оболочки,
свой процесс. Запуск команды внутри оболочки порождает
\textbf{подпроцесс}. Подпроцесс наследует состояние родителя
(текущий каталог, переменные), но это \textbf{копия}: изменения
в подпроцессе не влияют на родителя.

\begin{minted}[bgcolor=gray!6]{bash}
cd ~/seminar1/source     # сменить каталог в текущей оболочке
bash                     # запустить вложенную оболочку: дочерний процесс
# внутри неё:
cd source                # сменится только каталог вложенной оболочки
touch tmp/x.txt          # файл создаётся в общей ФС и остаётся
exit                     # завершить вложенную оболочку
\end{minted}

После \cmd{exit} текущий каталог ``откатывается'' к значению
родительской оболочки, а созданный файл на месте --- файловая система
общая, а каталог --- состояние процесса.

\begin{tipbox}
Перейти в каталог \textbf{из скрипта} для всей текущей сессии
не получится: \cmd{cd} в скрипте выполняется в подпроцессе.
Поэтому для изменения окружения текущей сессии используют
\emph{подключение} файла (\cmd{source} или \cmd{.}),
см.\ раздел \ref{sec:bashrc}.
\end{tipbox}

\section{Переменные окружения}

\textbf{Переменная окружения} --- именованная строка, доступная
процессу. В оболочке переменная читается через \cmd{$NAME}:
\begin{minted}[bgcolor=gray!6]{bash}
echo $HOME        # домашний каталог
echo $USER        # текущий пользователь
echo $SHELL       # путь к оболочке по умолчанию
echo $PATH        # каталоги поиска программ
\end{minted}

\cmd{PATH} --- список каталогов, в которых оболочка ищет исполняемые
файлы при вводе команды. Именно благодаря \cmd{PATH} вы пишете
\cmd{ls}, а не \cmd{/usr/bin/ls}.

Создание переменной:
\begin{minted}[bgcolor=gray!6]{bash}
FOO=bar           # локальная переменная текущей оболочки
echo $FOO
\end{minted}

Локальная переменная \textbf{не видна} подпроцессам. Чтобы передать
её запускаемым программам, нужно \textbf{экспортировать}:
\begin{minted}[bgcolor=gray!6]{bash}
export FOO
# теперь дочерний процесс её видит:
bash -c 'echo $FOO'
\end{minted}

\cmd{export} переводит переменную в \emph{окружение} процесса,
которое наследуется дочерними. Список всех переменных окружения
выводит \cmd{env} или \cmd{printenv}; сравнить два окружения
удобно через \cmd{diff}.

\section{Файлы настройки оболочки}
\label{sec:bashrc}

Оболочка при запуске читает файлы настройки. Основные:

\begin{itemize}
  \item \fpath{/etc/profile} и \fpath{/etc/bash.bashrc} ---
    \emph{системные}, одинаковы для всех пользователей;
  \item \fpath{~/.profile} --- читается
    при \emph{входной} (login) сессии;
  \item \fpath{~/.bashrc} --- читается
    при \emph{интерактивной} сессии (открытие терминала).
\end{itemize}

На практике почти всё добавляют в \fpath{~/.bashrc}.
Например, чтобы переменная появлялась в каждой новой оболочке:
\begin{minted}[bgcolor=gray!6]{bash}
echo 'export COURSE="Operating Systems and Development Tools"' \
    >> ~/.bashrc
\end{minted}

Изменение вступит в силу в \textbf{новой} оболочке. Чтобы применить
сразу в текущей, подключите файл:
\begin{minted}[bgcolor=gray!6]{bash}
source ~/.bashrc     # или: . ~/.bashrc
\end{minted}

\cmd{source} выполняет команды \emph{в текущей} оболочке,
а не в подпроцессе --- именно поэтому переменная и остаётся.

\section{Типичные ошибки}
\begin{itemize}
  \item \emph{command not found} --- команда не найдена в \cmd{PATH}.
    Проверьте орфографию, установлен ли пакет; для полного пути
    используйте \cmd{which} / \cmd{type}.
  \item Изменение переменной ``не сохранилось'' в новый терминал ---
    вы создали локальную переменную и не экспортировали/не
    добавили в \fpath{.bashrc}.
  \item \emph{Permission denied} при запуске скрипта --- нет права
    на исполнение (см.\ \nameref{ch:05-permissions});
    можно запустить явно: \cmd{bash script.sh}.
  \item Скрипт ``меняет'' каталог/переменные, но после его запуска
    ничего не изменилось --- скрипт выполняется в подпроцессе;
    используйте \cmd{source}.
\end{itemize}

\section{Связь с остальным курсом}
Оболочка --- базис для всего остального: пакеты ставятся
командами (\nameref{ch:08-packages}), код пишется в редакторе
(\nameref{ch:09-vim}), сессии организуются в \nameref{ch:10-tmux},
а сценарии из команд собираются в
\emph{скрипты оболочки} (\nameref{ch:27-scripting-ide}).
Мониторинг процессов, запущенных оболочкой, ---
\nameref{ch:12-processes}.

\section{Задачи для самопроверки}
\begin{enumerate}
  \item Объясните разницу между встроенной командой и внешней
    программой; приведите по два примера.
  \item Создайте локальную переменную и покажите, что подпроцесс
    её не видит; затем экспортируйте и покажите, что видит.
  \item Разместите в \fpath{.bashrc} переменную, проверьте её в
    новой оболочке и объясните, почему \cmd{source} нужен,
    чтобы применить изменения сразу.
  \item Запустите вложенный \cmd{bash}, поменяйте в нём каталог
    и завершите его. Объясните, почему текущий каталог
    родительской оболочки не изменился, а созданный файл остался.
\end{enumerate}
